\section{Indexed algebraic data types}
\applabel{indexed adts}

In \subsecref{indexed data types} we summarised the changes to the
simply-typed development of symmetric data types. In this appendix, we
spell the theory out in full detail, for completeness.

To extend the theory to indexed data types $\carrier\Ctx \types
\TConst D$, we use endofunctors $\struct[\TConst D] :
\Fam_{\carrier\Ctx} \to \Fam_{\Ctx}$ over the category
$\Fam_{\carrier\Ctx}$ of families indexed by a set $\carrier\Ctx$ and
vertical maps between them. The data type $\TConst D$ denotes an
initial $\struct[\TConst D]$-algebra $\carrier\Ctx \types \roll :
\struct[\TConst D] \TConst D \to \TConst D$. As in the non-indexed
case, we equip those endofunctors with a symmetry strength relative to
an ambient $\G$-action $\Ctx = \pair{\carrier\Ctx}{(\act)}$ on the set
of indices. This strength will be a certain natural transformation
$\strength : \gast \struct[\TConst D] \X \to \struct[\TConst D](\gast
\X)$, where $(\gast-)$ is the functor part of the \emph{free
$\Ctx$-indexed $\G$-action} monad.  Initial algebras for such
endofunctors enjoy an indexed version of the parameterised initiality
and representation theorems, and stating and establishing those is our
main goal.

We'll use the indexed family $\mt : \Tree \types \NodeIn \mt$ as our
running example, but unlike \exampleref{node lookup}, we avoid the
redundant cases through symmetric programming. As our endofunctor, we
take:
\begin{gather*}
(\struct[\NodeIn]\X)_{\Leaf} \definedby \emptyset
\qquad
(\struct[\NodeIn]\X)_{\Node\,\ell\, n\,r} \definedby
  \Nil \Or \Lft\:: \X_{\ell} \Or \Rgt\:: \X_{r}
\\
(\fmap f)_{u}\Nil \definedby \Nil
\qquad
(\fmap f)_{u}(\Lft\:: x) \definedby \Lft\:: f_{\ell}x
\qquad
(\fmap f)_{u}(\Rgt\:: y) \definedby \Rgt\:: f_{r}y
\tag{where: $\mt : \Tree \types f : \X \to \Y$, $u = \Node\, \ell\, n\, r$}
\end{gather*}
The initial algebra has as carrier $\mt :
\Tree \types \NodeIn\mt$ from \exampleref{families example}.

We denote the multiplication
of the free $\Ctx$-indexed $\G$-action monad $P$ by:
\[
(\muast) : \gast(\gast \X) \to \gast\X
\qquad
\pair g{\pair h{x : \X_{\inv h \act \inv g \act \ctx}}} \definedby
\pair{g\mul h}{x : \X_{\inv*{g\mul h}\act\ctx}}
\]
It is the functor $(\gast)$ that generalises the notion of strength to
the indexed setting:
\begin{definition}\definitionlabel{indexed symmetry strength}
  Let $\Ctx$ be a $\G$-action and $F : \Fam_{\carrier\Ctx} \to
  \Fam_{\carrier\Ctx}$ be an endofunctor. A \emph{$\Ctx$-indexed
  $\G$-symmetry strength} is a natural transformation $\strength :
  \gast F\X \to F(\gast \X)$, satisfying:
  \begin{gather*}
  \strength_{\ctx}(\unit, u) = \fmap(\fun x\pair\unit x) u
  \qquad
  \fmap(\muast)
  \compose
    \strength(g, \strength(h, x)) =
  \strength(g\mul h, x)
  \end{gather*}
\end{definition}

\begin{example}%
  \examplelabel{path strength}
  We define an $\Sym2$-symmetry strength for $\mt : \Tree \types \NodeIn\mt$.
  Only the $\swap$ case is interesting, as the first strength axiom forces the
  definition for $\stet$:
  \begin{gather*}
  \strength[\NodeIn]_{\parent{\NodeV \ell n r}}
  \pair{\text{$\swap$}}{\Nil : \NodeIn {\NodeV{\swap\act r}n{\swap\act \ell}}}
  \definedby \Nil : \NodeIn \parent{\NodeV{\ell}n{r}}
  \\
  \strength[\NodeIn]_{\parent{\NodeV \ell n r}}\pair\swap{\Lft \:: (x : \X_{\swap\act r})} \definedby
  \Rgt \:: \pair\swap{x : \X_{\swap \act r}} \in (\struct[\NodeIn](\gast\X))_{\parent{\NodeV \ell nr}}
  \\
  \strength[\NodeIn]_{\parent{\NodeV \ell n r}}\pair\swap{\Rgt \:: (x : \X_{\swap\act \ell})} \definedby
  \Lft \:: \pair\swap{x : \X_{\swap \act \ell}} \in (\struct[\NodeIn](\gast\X))_{\parent{\NodeV \ell nr}}
  \end{gather*}
  To prove the second strength axiom, only the case $g \definedby h
  \definedby \swap$ is interesting though straightforward:
  \begin{align*}
    \fmap(\muast)&\compose\strength(\swap, \strength(\swap, \Nil))
    = \fmap(\muast)\compose\strength(\swap, \Nil)
    = \fmap(\muast)\Nil
    = \Nil
    \\&= \strength(\swap\mul\swap, \Nil)
    \\
    \fmap(\muast)&\compose\strength(\swap, \strength(\swap, \Lft\:: x))
    = \fmap(\muast)\compose\strength(\swap, \Rgt\::\pair\swap x)
    \\&= \fmap(\muast)(\Lft\::\pair{\swap\mul\swap} x)
    = \Lft \:: \pair\stet x
    = \strength(\swap\mul\swap, \Lft \::x)
  \end{align*}
  and symmetrically for $\Rgt::x$. We cannot avoid this
  symmetric argument using wlog as we are setting up the symmetric
  action by defining the strength. Later uses of $\wlg{}$
  will exploit this work.
\end{example}

The appropriate category-theoretic concept that abstracts this notion
of strength is a \emph{monad-functor distributive law}, and we will
use this perspective in our category theoretic development in
\appref{cat dev}. As may be clear, when $\terminal \types X$ is a
single-fibred family, $\terminal \types \gast X$ is the single-fibred
family $\carrier G \times X$, subsuming the simply-typed case.

The following lemma justifies checking only the interesting cases, as
we did in \exampleref{path strength}:
\begin{lemma}\lemmalabel{indexed check only non-trivial}
  Let $\strength$ be the structure of an indexed $\G$-symmetric
  strength, satisfying the first axiom. Then the second axiom holds
  when either $g = \unit$ or $h = \unit$.
\end{lemma}

Our main initiality theorem uses an indexed notion of equivariant
algebras (cf.~\definitionref{equivariant F-algebra}):
\begin{definition}
  Let $\Ctx$ be a $\G$-action, $F : \Fam_{\carrier\Ctx} \to
  \Fam_{\carrier\Ctx}$ an endofunctor equipped with a symmetry
  strength. For every indexed action $\A$, the \emph{lifted action}
  is given by:
  \[
  \carrier{F_{\strength}\A} \definedby F\carrier\A
  \qquad
  g \iact\ctx u \definedby \fmap (\act) \strength_{g\act\ctx}\pair gt
  \]
  An \emph{equivariant $F$-algebra} $\pair\A f$ for an endofunctor
  equipped with a symmetry strength, consists of an indexed action
  $\A$ equipped with an equivariant algebra structure $\Ctx \types f :
  F_{\strength}\A \equito \A$. An equivariant $F$-algebra homomorphism is
  then a family map between the carrier families of two such algebras
  that is simultaneously equivariant and $F$-algebra homomorphic.
\end{definition}

We generalise the Equivariant Initiality \theoremref{representation for algebras}
to indexed families:
\begin{theorem}[indexed equivariant initiality]%
  \theoremlabel{indexed representation}
  Let $\roll : F\fix F \to \fix F$ be an initial algebra for a
  functor $F$ over $\Ctx$-indexed families with a $\G$-symmetry strength
  $\strength : \gast
  F\X \to F(\gast\X)$. There is a unique map $(\act) : \gast \fix F \to \fix F$
  satisfying \(
  g \act (\roll u)
  =
  \roll (\fmap(\act)(\strength(g, u)))
  \).
  It makes $\pair{\fix f}{\roll}$
  into an initial equivariant $F$-algebra: for every
  equivariant $F$-algebra $\pair\A{f : F_{\strength}\A \equito \A}$,
  the
  $F$-homomorphism into $\carrier\A$
   is equivariant: $\fold f :
  \fix F \equito \A$.
\end{theorem}

Like its simply-typed variant, the proof of this theorem uses an
appropriate notion of parameterised initiality originally proposed by
\citeauthor{bird-paterson:generalised-folds-for-nested-datatypes}~\cite{%
  bird-paterson:generalised-folds-for-nested-datatypes}.

\begin{example}
  The action induced by the Indexed Equivariant Initiality
  \theoremref{indexed representation} satisfies:
  \begin{align*}
  \swap\act(\Lft \:: \Rgt \:: \Nil) &=
  \roll\compose\fmap(\act)\strength(\swap, \Lft \:: (\Rgt \:: \Nil)) =
  \roll\compose\fmap(\act)(\Rgt \:: \pair\swap{\Rgt \:: \Nil})
  \\&=
  \Rgt\::(\swap \act (\Rgt \:: \Nil))=
  \Rgt\::\Lft \:: (\swap \act \Nil)=
  \Rgt\::\Lft \:: \Nil
  \end{align*}
  as we may expect from the example at the beginning of \subsecref{indexed actions}.
\end{example}

\begin{example}%
  \examplelabel{wlog node lookup}%
  Let's implement the dependently-typed node lookup function and
  simultaneously show its equivariance $\lookup :
  \DPair\mt{\Tree}{\NodeIn\mt} \equito \naturals$. We'll
  construct an indexed equivariant map $\mt : \Tree \types
  \lookup_{\mt} : \NodeIn\mt \equito \naturals$. By the Indexed
  Equivariant Initiality \theoremref{indexed representation}, it
  suffices to equip $\naturals$ with an equivariant
  $\struct[\NodeIn]$-algebra structure $\lookupElim : \struct[\NodeIn]
  \naturals \equito \naturals$. Informally, \WLOG{}, the path to the
  node is either empty or starts with the $\Lft$eft subtree,
  eliminating one symmetric case. Formally, define a
  canoniser by
  \(
  \canoniser\Nil \definedby \canoniser(\Lft\:: n) \definedby \stet
  \) and \(
  \canoniser(\Rgt\:: n) \definedby \swap
  \).
  Its canonical elements are $\Nil$ and $\Lft\::n$.
  This canoniser is equivariant, but not strictly:
  \begin{align*}
  \canoniser(\swap \act \Nil) &=
  \canoniser\Nil = \stet \stabeq{\Nil} \swap = \stet\mul\inv\swap =
  \swap\act \canoniser\,\Nil \\
  \canoniser(\swap \act (\Lft \:: n)) &=
  \canoniser(\Rgt \:: n) =
  \swap = \stet \mul \inv\swap = \swap\act\canoniser(\Lft\:: n)
  \end{align*}
  and symmetrically for $\Rgt \:: n$. Again, we cannot avoid this
  symmetry with wlog since we are setting up the symmetric argument.

  Let's define the core function for
  the eliminator $\mt : \Tree \types \lookupCore :
  \Canon\canoniser \to \naturals$:
  \[
  \lookupCore_{\parent{\Node\, \ell\,n\,  r}} \Nil \definedby n
  \qquad
  \lookupCore_{\parent{\Node\,\ell\,n\,  r}} (\Lft \:: p) \definedby p
  \]
  and we do not need to consider the symmetric case.
  The stabiliser subgroup for $\naturals$ is the whole of $\Sym2$
  since its action is trivial. Consequently, every core function is
  stable, and we don't need to work out the stabilisers for canonical
  elements. For completeness, here they are.  The stabiliser subgroup
  for $\Nil$ is the whole of $\Sym2$, and the stabiliser subgroup for
  $\Lft\::p$ is the trivial subgroup $\set\stet$.

  We use the
  Indexed Representation \theoremref{indexed-wlog-construction} to
  show the extended function is equivariant: $\lookupElim \definedby
  \wlg(c,f) : \struct[\NodeIn]\naturals \equito \naturals$.
  \end{example}

As in the simply-typed case (cf.~\subsecref{wlog-traversals}), we
reuse the same canoniser for multiple traversals of the same symmetric
indexed data type, and streamline the steps in \exampleref{wlog node lookup}.

\begin{definition}
  Let $\Ctx$ be a $\G$-action, and $F$ be a $\Ctx$-indexed
  $\G$-strong.  An \emph{indexed canoniser} for $F$ is a natural
  transformation $\carrier\Ctx \types \canoniser : F \to
  \carrier\G$. Each such canoniser induces the \emph{canonical
  elements} functor $\Canon\canoniser : \Action\G_{\Ctx} \to
  \Fam_{\carrier\Ctx}$ given by:
  \[
  ((\Canon\canoniser)\A)_{\ctx} \definedby
  (\Canon[F\A]\canoniser_{\carrier\A})_{\ctx} \definedby
  \set{
    g \iact[F\A]\dtx u \suchthat
    \begin{array}{@{}l@{}}
      g = \canoniser_{\carrier\A; \dtx}u, g \act \dtx = \ctx,\\
      u \in (\carrier F\carrier\A)_{\dtx}
    \end{array}
  }
  \qquad
  (\Canon\canoniser)h : \fun u (\carrier Fh)u
  \]
\end{definition}

An indexed canoniser is uniquely determined by its component at the
terminal family $\carrier\Ctx \types \terminal$.

\begin{example}
  The indexed $\Sym 2$-strong functor $\struct[\NodeIn]$ has the
  canoniser (cf.\exampleref{goLeft canoniser})
  \[
  \goLeft_{\Node n\,\ell\,r}\,\Nil \definedby \stet \defines \goLeft_{\Node \ell\,n\,r} (\Lft\:: p)
  \quad
  \goLeft_{\Node \ell\,n\,r} (\Rgt\:: p) \definedby \swap
  \]
  The canonical elements are:
  \[
  (\Canon[\A]\goLeft)_{\Leaf} \definedby \emptyset
  \qquad
  (\Canon[\A]\goLeft)_{\Node \ell\,n\,r} \definedby
  \set{\Nil, \Lft\:: a \suchthat a \in \carrier\A_{\Node \ell\,n\,r}}
  \tag*{\qedhere}
  \]
\end{example}

Continuing in tandem with the development of \subsecref{indexed data types},
we define the data we use to traverse an indexed symmetric data type:

\begin{definition}\definitionlabel{indexed core algebras}
  Let $F$ be a $\Ctx$-indexed $\G$-strong functor and $\canoniser : \carrier F \to
  \carrier\G$ a canoniser for it. An \emph{indexed $\canoniser$-core
  $F$-algebra} $\pair\A f$ consists of a $\Ctx$-indexed $\G$-action $\A$ equipped
  with a vertical map $\carrier\Ctx \types f : (\Canon\canoniser)\A \to \carrier\A$. It is
  \emph{equivariant} when $f$ is equivariant $\Ctx \types f :
  \parent{(\Canon\canoniser)\A \subset F\A} \equito \A$. A
  \emph{homomorphism} of equivariant indexed $\canoniser$-core $F$-algebras $h
  : \pair\A f \to \pair \B g$ is an equivariant map $\Ctx \types h : \A \equito
  \B$ satisfying: $h (f\, u) = g\, ((\Canon\canoniser) h\,u)$.
\end{definition}

\begin{example}\examplelabel{dependently typed lookup as core algebra}
  We have the following equivariant $\goLeft$-core $\struct[{\NodeIn}]$-algebra:
  \begin{gather*}
  \lookupCore : (\Canon\goLeft)(\Triv{\naturals}) \to
  \carrier{\Triv{\naturals}} = \naturals
  \\
  \lookupCore_{\Node \ell\,n\,r}
  \Nil \definedby n
  \qquad
  \lookupCore_{\Node \ell\,n\,r}
  \parent{
    \Lft\:: (k \in \naturals)
  } \definedby k
  \tag*{\qedhere}
  \end{gather*}
\end{example}

Next, chain together the traversal and \WLOG{} constructions:
\[
\carrier\Ctx\types
\parent{\Canon[F\A]\canoniser_{\carrier\A} = (\Canon\canoniser)\A \xto f \carrier\A}
\mapsto
\parent{F\carrier \A \xto{\wlg(\canoniser_{\carrier\A}, f)} \carrier\A}
\mapsto
\parent{\flg[\A](\canoniser, f) :
  \fix\carrier F \xto{\fold \wlg(\canoniser_{\carrier\A}, f)} \carrier \A}
\]
And similarly chain \theoremref{indexed-wlog-construction} and the
Indexed Equivariant Initiality \theoremref{indexed representation}:


\begin{theorem}[indexed core initiality]
  \theoremlabel{indexed representation for core algebras}
  Let $F$ be a $\Ctx$-indexed $\G$-strong functor, $\carrier\Ctx \types
  \roll : \carrier F\fix
  \carrier F \to \fix \carrier F$ an initial algebra of $\carrier F$,
  and $\canoniser$ a canoniser for $F$. Then restricting $\roll$ to canonical elements
  makes $\fix F$ into the initial indexed $\canoniser$-core $F$-algebra
  $\roll : (\Canon\canoniser)\fix F \equito \fix F$:
  for any equivariant indexed $\canoniser$-core $F$-algebra ${f :
    (\Canon\canoniser)\A \equito A}$, the unique equivariant
  indexed $\canoniser$-core $F$-algebra homomorphism
  into $f$ is:\/ $\flg(\canoniser, f) : \fix F \equito \A$.
\end{theorem}

This theorem implies, by \exampleref{dependently typed lookup as core
  algebra}, that we have an equivariant function:
\[
t : \Tree \types \lookup \definedby
\flg(\goLeft, \lookupCore) : \text{$\NodeIn t \equito \Triv \naturals$}
% For some reason this empty \mathbf is needed, otherwise the font in
% \flg at the beginning on the line preceding this comment is typeset
% in blackboard bold
\mathbf{}
\]
